%% ma: L10-B12-0007
%% lop: 10
%% chuong: 5
%% bai: 12
%% dang: 10.12.8
%% loai: TL
%% muc: TH
%% dap_an: "-"
%% nguon: "chuyen-de-cac-so-dac-trung-cua-mau-so-lieu-khong-ghep-nhom-toan-10-chuong-trinh-moi.pdf (GV. Phan Nhật Linh) (P01), Bài 01 Bài tập tự luận 7"
%% kien_thuc: "Lớp 10 Bài 12 (số gần đúng, sai số tuyệt đối, độ chính xác, sai số tương đối, quy tròn số)"
%% trang_thai: chinh-thuc
%% ghi_chu: "Tự duyệt theo yêu cầu giáo viên 22/10/2026."
\begin{ex}
Minh tính diện tích của một hình tròn có bán kính $r=4$ cm bằng công thức $S=3{,}145\cdot4^2=50{,}32\ (\text{cm}^2)$. Biết rằng $3{,}14<\pi<3{,}15$. Hãy ước lượng độ chính xác của $S$.
\loigiai{Diện tích đúng $\overline S=16\pi$ thỏa mãn $3{,}14\cdot16<\overline S<3{,}15\cdot16$, tức là $50{,}24<\overline S<50{,}4$.

Suy ra $-0{,}08<\overline S-50{,}32<0{,}08$, tức là $\left|\overline S-50{,}32\right|<0{,}08$.

Vậy kết quả của Minh có độ chính xác $0{,}08$.}
\end{ex}
